\begin{claim}{D.48}{Proved}{The two sections of a punctured-torus order-three action}
Let $a$ have order three in the full once-punctured-torus
loop-outer group of $F=\widehat F(x,y)$, and let $E_a$
be its centre-free extension.
There are exactly two $F$-conjugacy classes of order-three
sections in $E_a$.  Each such section $S$ satisfies
\[
                         H^1(C_F(S),\F_3)=0.             \tag{D.48a}
\]
The normalizer in $E_a$ of the torus cusp
$L=\overline{\langle[x,y]\rangle}$ is torsion-free.
If $v$ is any element of this normalizer projecting
to the prescribed generator $a$, and
\[
             v[x,y]v^{-1}=[x,y]^\lambda,\qquad
                            v^3=[x,y]^\kappa ,
\]
then its $3$-adic coordinate satisfies
$\kappa_3\in\Z_3^\times$.
\end{claim}
\begin{proof}
\emph{The actual action and its abelian quotient.}
D.27 gives an actual order-three representative
$\alpha$ of $a$, so
$E_a=F\rtimes_\alpha\langle s\rangle$.
Let $A$ be its matrix on $H_1(F,\Z_3)$.
It has order exactly three by D.33.
Over $\Q_3$ the polynomial $X^2+X+1$ is irreducible:
a primitive cube root would give a square root of
$-3$, whose $3$-adic valuation is odd.
Since $A^3=I$ and $A\ne I$, it follows that
\[
             A^2+A+I=0,\qquad
             \det(A-I)=3,\qquad
             \operatorname{coker}(A-I)\cong C_3.
                                                        \tag{D.48b}
\]
The last equality follows by Smith reduction over
$\Z_3$, since the determinant has valuation one.
Abelianizing the actual semidirect product gives
the exact maximal pro-$3$ abelian quotient
\[
                     E_a^{\rm ab,(3)}\cong C_3^2.        \tag{D.48c}
\]
One factor is (D.48b), the other is the specified
section quotient.

The reduction of $A$ modulo $3$ is a nontrivial
unipotent Jordan block.  Indeed, if $A=I+3B$,
then $A^3=I$ implies
$B(I+3B+3B^2)=0$ after division by $9$.
The second factor is invertible, giving $B=0$,
contrary to $A\ne I$.
The induced action on
$M=H^1(F,\F_3)$ is therefore also one nontrivial
Jordan block of dimension two.
Writing $D$ for the action matrix on this dual
space, the exact cyclic resolution has successive
operators
\[
              D-I,\quad 1+D+D^2=0,\quad D-I,\quad0,\ldots,
 \qquad \operatorname{rank}(D-I)=1.
                                                        \tag{D.48d}
\]
Consequently $H^j(C_3,M)$ has dimension one for
every $j\geq0$.

\emph{Bockstein injectivity and periodicity.}
Let $u\in H^1(E_a,\F_3)$ be the quotient character,
and let $\chi$ be a nonzero invariant character of
$F$, extended to $E_a$ by $\chi(s)=0$.
They form a basis of $H^1(E_a,\F_3)$ by (D.48c).
For the coefficient sequence
\[
              0\longrightarrow\F_3
                 \xrightarrow{\;3\;}\Z/9
                 \longrightarrow\F_3\longrightarrow0
\]
write $\mathfrak b$ for the Bockstein.
Every homomorphism $E_a\to\Z/9$ factors through
(D.48c), so its values are in $3(\Z/9)$ and its
reduction modulo $3$ is zero.
The coefficient long exact sequence therefore makes
\[
     \mathfrak b:H^1(E_a,\F_3)\longrightarrow
                              H^2(E_a,\F_3)
                  \quad\hbox{injective}.                \tag{D.48e}
\]
Put $v=\mathfrak b(u)$ and $w=\mathfrak b(\chi)$.
The notation $v$ in this cohomology calculation
denotes a class, not the normalizer element in the
statement.
The class $v$ is the inflation of the usual degree-two
period class of $C_3$; it is nonzero on the section
$\langle s\rangle$.
The classes $v,w$ are linearly independent by
(D.48e).

The Hochschild--Serre spectral sequence has only
rows zero and one.  By (D.48d), at each positive
total degree its $E_2$ page has two
one-dimensional terms.  Thus its degree-two
abutment has dimension at most two.
The two independent Bockstein classes make its
dimension exactly two.
The only possible differential is
$d_2:E_2^{j,1}\to E_2^{j+2,0}$.
At $j=0$ it vanishes because $\chi$ extends to
$E_a$, and at $j=1$ it vanishes because the
one-dimensional term $E_2^{1,1}$ must survive
to obtain dimension two in degree two.

Multiplication by the period class $v$ identifies
each cyclic-resolution term with the term two
places later.
On the row with coefficient $M$ this also holds
starting at $j=0$, since the norm in (D.48d) is
zero, so $H^0(C_3,M)\to H^2(C_3,M)$ is the identity
on $\ker(D-I)$.
On the row with trivial coefficients it is the
usual period isomorphism.
The multiplicative spectral sequence therefore
makes every remaining $d_2$ a multiple of one of
the two already vanishing differentials.
It collapses, and multiplication by $v$ is an
isomorphism on its two associated graded pieces
at every positive total degree.  Hence
\[
 \begin{gathered}
        \dim H^n(E_a,\F_3)=2\quad(n\geq1),\\
        v:H^n(E_a,\F_3)\xrightarrow{\ \sim\ }
                             H^{n+2}(E_a,\F_3)
                       \quad(n\geq1),\\
                         wv^m\ne0\quad(m\geq0).
 \end{gathered}                                         \tag{D.48f}
\]

\emph{Exactly two full section classes.}
Every finite subgroup of $E_a$ injects into its
quotient $C_3$, because $F$ is torsion-free.
For an order-three section $S$ the normalizer is
the exact direct product
\[
                         N_{E_a}(S)=C_F(S)\times S.      \tag{D.48g}
\]
Conjugation by a normalizer element induces the
identity on $S$: projection to the abelian quotient
$C_3$ identifies every element of $S$ with its
label, so that label cannot change under
conjugation.
Multiplying by the section element with the same
quotient label then leaves precisely $C_F(S)$.

The full-profinite Brown restriction theorem
\cite[Theorem~2 and \S1.1]{A-KochloukovaZalesskii2007},
applied with the open subgroup $F$ of
$3$-cohomological dimension one, identifies
sufficiently high cohomology with the product
over (D.48g).
Since $C_F(S)$ is a closed subgroup of $F$,
its $3$-cohomological dimension is at most one.
Writing $k$ for the number of section classes,
the dimensions in (D.48f) therefore give
\[
                  k+\sum_{[S]}\dim H^1(C_F(S),\F_3)=2.
                                                        \tag{D.48h}
\]
In particular the class set is finite.

It cannot consist of a single class.
Otherwise take that class to be the chosen
$S=\langle s\rangle$.
On its normalizer (D.48g), the character $\chi$
is pulled back from $C_F(S)$, since $\chi(s)=0$.
Its Bockstein therefore restricts to zero:
\[
       w|_{N_{E_a}(S)}
        =\operatorname{pr}_{C_F(S)}^*
               \mathfrak b(\chi|_{C_F(S)})=0,
                                                        \tag{D.48i}
\]
because $H^2(C_F(S),\F_3)=0$.
Every nonzero class $wv^m$ from (D.48f) would
then have zero Brown restriction, a contradiction
for large $m$.
Thus $k\geq2$; (D.48h) makes $k=2$ and makes
both summands zero, proving (D.48a).
As in D.46, $E_a$-conjugacy equals $F$-conjugacy:
for any section $S$ choose its unique element
with the same quotient label as a given
conjugator and multiply the conjugator by its
inverse.

\emph{The peripheral normalizer.}
We spell out the boundary consequence of the
vanishing in (D.48a).
A closed subgroup $C\leq F$ with
$H^1(C,\F_3)=0$ cannot contain a conjugate of
the full $3$-part
$\langle[x,y]^{e_3}\rangle\cong\Z_3$
of the torus cusp.
First its supernatural index in $F$ must have
infinite $3$-part.  Otherwise choose an open
$U_0\geq C$ whose index has maximal $3$-valuation.
For every open $V$ with $C\leq V\leq U_0$,
the index $[U_0:V]$ is prime to $3$.
Transfer makes
$H^1(U_0,\F_3)\to H^1(V,\F_3)$ injective.
The restriction-colimit comparison of D.32
would then inject the nonzero group
$H^1(U_0,\F_3)$ into $H^1(C,\F_3)$,
a contradiction.

Choose an open $U\geq C$ with $3\mid[F:U]$.
Its corresponding finite cover of the punctured
torus has a lift of the selected cusp of degree
prime to $3$, because its completed boundary
subgroup contains the selected full $\Z_3$.
The exact discrete boundary comparison uses
the dense discrete representative for the
profinite conjugator, as in D.32 and D.33.
There are at least two boundary lifts: if there
were only one, its degree would be $[F:U]$.
The selected boundary class is nonzero in
$H_1(U,\F_3)$, since the only relation among
the boundary classes is their sum.
A boundary-dual character restricts nontrivially
to the selected $\Z_3$ in $C$; the scalar
comparison is the inverse of its covering
degree modulo $3$.
This again contradicts $H^1(C,\F_3)=0$ and proves
the assertion.

For $L=\overline{\langle[x,y]\rangle}$, the
peripheral normalizer has the exact sequence
\[
           1\longrightarrow L\longrightarrow N_{E_a}(L)
                    \longrightarrow C_3\longrightarrow1.
                                                        \tag{D.48j}
\]
Surjectivity follows from preservation of the cusp
class, and its kernel is $L$ by A.7.
The multiplier satisfies $\lambda^3=1$, and
$\lambda_3=1$, since $\Z_3^\times$ has no
nontrivial root of unity of order three.
An order-three section in (D.48j) would therefore
centralize the full cusp $\Z_3$.
Its fixed subgroup in $F$ would contradict the
boundary consequence of (D.48a).
Every torsion subgroup of (D.48j) injects into
$C_3$, so this proves that the normalizer is
torsion-free.

Finally take the element called $v$ in the
statement, with $v^3=[x,y]^\kappa$.
Commuting $v$ with its own cube gives
$(\lambda-1)\kappa=0$.
Where $\lambda_p\ne1$, the coordinate $\kappa_p$
and the norm $1+\lambda_p+\lambda_p^2$ are zero.
Where $\lambda_p=1$, that norm is $3$.
If $\kappa_3$ were divisible by $3$, the exact
norm equation
\[
                (1+\lambda+\lambda^2)z=-\kappa
                               \quad\text{in }\widehat\Z
                                                        \tag{D.48k}
\]
would consequently have a solution: divide by
$3$ at the positive coordinates, and use zero
at the other coordinates.
The element $v[x,y]^z$ would then have order
three and project to $a$, contradicting the
torsion-free conclusion.
Thus $\kappa_3$ is a unit, as asserted.

The program
\texttt{code/part\_c/torus\_order3\_sections\_d48.py}
checks the integral and mod-$3$ matrix calculations,
the abelian character and Bockstein input in its
finite semidirect quotient, and the cyclic
resolution operators.
The full restriction and finite-cover arguments
above establish the section and normalizer
assertions.  They do not assert global triviality
of the order-three multiplier or geometric
conjugacy.
\end{proof}
